% =====================================================================
%  PETase-ML — Two-stage, leakage-audited stability framework
%  Nature Scientific Reports submission draft — Overleaf-ready.
%
%  HOW TO USE ON OVERLEAF
%   1. Create a new project, paste this whole file as main.tex.
%   2. Upload the figure file  precision_vs_threshold.pdf  (from
%      paper_figures/) into the project root, next to this .tex.
%   3. Compile with pdfLaTeX.
%
%  Everything highlighted in YELLOW  [TO ADD: ...]  is a real gap that an
%  author must fill or confirm before submission. Confirmed numbers are
%  taken from the study result files and are already in place.
%
%  SCOPE / VERSION CONTROL (editorial — not printed):
%   * Temperature-focused study; pH-optimum analysis excluded (future work);
%     EpHod mentioned only as future work.
%   * General protein-stability framework evaluated for potential application
%     to PET-degrading enzyme engineering (NOT PETase-specific training data).
%   * No experimental / wet-lab / industrial / environmental claims.
%   * Older 19,071-mutation results (acc ~79.7%, r ~0.764, MAE ~0.92) are
%     superseded by the 1,001,888-row two-stage results and must NOT appear.
% =====================================================================
\documentclass[11pt]{article}

\usepackage[margin=1in]{geometry}
\usepackage{graphicx}
\usepackage{booktabs}
\usepackage{array}
\usepackage{amsmath,amssymb}
\usepackage[table]{xcolor}
\usepackage{authblk}
\usepackage{caption}
\usepackage{lineno}
\usepackage[colorlinks=true,allcolors=blue]{hyperref}

% ---- highlighted "to add" markers -----------------------------------
\definecolor{filleryellow}{RGB}{255,238,140}
\newcommand{\filler}[1]{\colorbox{filleryellow}{\bfseries [TO ADD: #1]}}
\newcommand{\fillerblock}[1]{%
  \par\smallskip\noindent
  \colorbox{filleryellow}{\parbox{\dimexpr\linewidth-2\fboxsep\relax}{%
    \small\textbf{[TO ADD]}\ #1}}\par\smallskip}

% ---- shorthand for units / symbols ----------------------------------
\newcommand{\ddg}{\ensuremath{\Delta\Delta G}}
\newcommand{\Tm}{\ensuremath{T_\mathrm{m}}}
\newcommand{\dtm}{\ensuremath{\Delta T_\mathrm{m}}}
\newcommand{\degC}{\ensuremath{^{\circ}}C}

\captionsetup{labelfont=bf,font=small}
\setlength{\parskip}{0.4em}

% =====================================================================
\title{\bfseries A Two-Stage, Leakage-Audited Machine-Learning Framework for
Temperature- and Stability-Aware Screening of PET-Degrading Enzymes}

\author[1]{Ayush Iyer\thanks{Corresponding author: \texttt{iyer.ayush31@gmail.com}}}
\author[2]{James Ponzio}
\author[3]{Abhinav Iyer}
\affil[1]{\filler{affiliation --- institution and country for Ayush Iyer}}
\affil[2]{\filler{affiliation --- institution and country for James Ponzio}}
\affil[3]{\filler{affiliation --- institution and country for Abhinav Iyer}}
\date{}

% =====================================================================
\begin{document}
\maketitle
\linenumbers

\noindent\textit{Title placeholder note:} \filler{confirm final title, $\leq$20 words, no undefined abbreviations}

% =====================================================================
\section*{Abstract}
Plastic pollution is a defining environmental challenge, and enzymatic depolymerization of
polyethylene terephthalate (PET) offers a biological route to recycling; however, candidate
enzymes must remain stable under demanding conditions, and computational screening is
frequently inflated by data leakage. We built a general protein-stability framework and
evaluated it for potential application to PET-degrading enzyme engineering. We assembled
1{,}001{,}888 mutation and stability records from seven public sources and routed them by
measurement type into a two-stage design: Stage~1 screens whole-protein thermostability
(melting temperature, \Tm) and Stage~2 optimizes mutation-level folding stability (\ddg).
Training records were audited for exact and sequence-homology leakage against the independent
benchmarks. On the independent BRENDA benchmark, Stage~1 reached a ROC AUC of 0.732 (95\% CI
0.708--0.755). On the independent S669 benchmark, Stage~2 reached a ROC AUC of 0.669 (95\% CI
0.62--0.72) and an ensemble Pearson correlation of 0.390 (95\% CI 0.32--0.46). Both stages
showed a precision--recall trade-off in which precision rose and recall fell under stricter
thresholds. These results describe a moderate, leakage-controlled screen-then-optimize
framework suited to prioritizing candidates for experimental testing, not a validated
industrial solution.
\fillerblock{confirm final word count $\leq$200; ``ROC AUC'' denotes receiver-operating-characteristic AUC (not precision--recall AUC); internal validation is not presented as independent testing.}

\medskip
\noindent\textbf{Keywords:} protein stability; machine learning; enzyme engineering;
thermostability (\Tm); mutation-effect (\ddg) prediction; PET-degrading enzymes.

% =====================================================================
\section{Introduction}
Polyethylene terephthalate (PET) is among the most heavily produced synthetic polymers, and its
accumulation as waste is a major contributor to global plastic pollution
\filler{cite primary source; give exact figure and state production/waste/municipal/PET}.
Biological recycling, in which enzymes depolymerize PET into recoverable monomers, is an
attractive complement to mechanical and chemical recycling \filler{ref}.

Enzymatic PET degradation has advanced through the discovery and engineering of PET hydrolases
such as PETase and related cutinases \filler{ref}. Thermostability matters because reactions run
faster near PET's glass-transition temperature, but it is not the sole determinant of catalytic
performance: activity, substrate binding, reaction conditions, polymer crystallinity, and pH
also matter \filler{ref}.

Computational stability prediction addresses two related but non-interchangeable questions.
Whole-protein screening asks whether a given enzyme is thermostable, summarized by its melting
temperature (\Tm). Mutation-level optimization asks how a specific substitution changes folding
stability, expressed as the change in unfolding free energy (\ddg). Conflating the two leads to
misleading claims about what a model can do.

A recurring problem in this literature is data leakage. Exact-match leakage (identical sequences
or mutations shared between training and test) and, more subtly, sequence-homology leakage
(shared homologous background) can inflate benchmark performance \filler{ref}. We do not claim
that all prior models contain leakage; rather, we control for it explicitly and report the
effect.

Because whole-protein selection and mutation-level optimization are usually treated separately,
this study evaluates a single, integrated two-stage screen-then-optimize framework---a general
protein-stability method assessed for potential application to PET-degrading enzyme
engineering---on leakage-audited, independent BRENDA (temperature) and S669 (\ddg) benchmarks,
and reports its performance without overstating readiness.

% =====================================================================
\section{Results}

\subsection{Dataset composition and leakage auditing}
After harmonization, the multi-source staging table comprised 1{,}001{,}888 mutation and
stability records drawn from seven public sources (Table~\ref{tab:data}, Panel~A). Because the
framework treats temperature stability and folding stability as distinct problems, records were
routed by measurement type (Table~\ref{tab:data}, Panel~B): 29{,}654 melting-temperature (\Tm)
records were assigned to Stage~1 and 508{,}693 folding free-energy change (\ddg) records to
Stage~2. Abundance-proxy (457{,}943) and \dtm{} (5{,}598) records were retained in the table but
were not used for model training.

To prevent optimistic bias, every candidate training record was audited against the independent
test proteins at two levels: exact-match (identical sequence and mutation) and sequence-homology
(shared homologous background). Auditing removed approximately 132{,}000 records, of which
131{,}479 were flagged by homology alone and would have been invisible to exact-match filtering.
This indicates that homology-level leakage, rather than duplicate records, is the dominant
contamination risk in this setting.

\subsection{Stage 1: temperature-stability performance (BRENDA)}
\textit{Internal validation (development only).} During development, the Stage~1 temperature
model achieved an internal held-out RMSE of approximately 6.5~\degC{} and a validation Pearson
correlation of approximately 0.80. These figures describe performance on data drawn from the
training distribution and are reported solely to characterize model fitting; they are not
evidence of generalization and are kept separate from the independent results below.

\textit{Independent testing (BRENDA benchmark).} Independent evaluation used a curated BRENDA
whole-protein temperature-stability benchmark. During curation, 412 records were removed to yield
2{,}034 proteins; a further 471 ambiguous cases were excluded, leaving 1{,}563 proteins (979
positive, 584 negative). On this independent set the classifier reached a ROC AUC of 0.732 (95\%
CI 0.708--0.755). Evaluated as a regressor against measured \Tm, it obtained a Pearson
correlation of 0.54 (95\% CI 0.50--0.57), rising to 0.60 on the subset with directly measured
true-\Tm{} values (Table~\ref{tab:stage1}).

\textit{Threshold behaviour.} Classification performance depended strongly on the decision
threshold (Fig.~\ref{fig:precision}a). As the temperature cutoff was made more stringent,
precision increased while recall fell---from 0.65 (recall 0.98) at the most permissive cutoff to
0.93 (recall 0.39) and 0.98 (recall 0.32) at the strictest cutoffs (Table~\ref{tab:stage1},
Panel~A). We emphasize that precision at a strict threshold is not an experimental success rate:
the 0.93 precision corresponds to only 39\% recall, so the strict-threshold regime identifies a
small, high-confidence set of stabilizing candidates while missing most true positives.

\subsection{Stage 2: \texorpdfstring{\ddg}{ddG} stabilization performance (S669)}
Independent folding-stability performance was assessed on the S669 benchmark of 669 single
mutations (168 stabilizing, 501 non-stabilizing reference cases). Across 2{,}000 bootstrap
resamples, the gradient-boosted ensemble achieved an ensemble Pearson correlation of 0.390 (95\%
CI 0.32--0.46), RMSE 1.509~kcal\,mol$^{-1}$, and a ROC AUC of 0.669 (95\% CI 0.62--0.72) for
stabilizing/non-stabilizing classification (Table~\ref{tab:stage2}, Panels~A--B). Among
individual models, Extra Trees---not the ensemble---produced the strongest single-model
regression metrics (Pearson 0.436; Table~\ref{tab:stage2}, Panel~A).

As in Stage~1, precision and recall traded off with threshold stringency
(Fig.~\ref{fig:precision}b; Table~\ref{tab:stage2}, Panels~B--C). At the strict threshold the
ensemble reached precision 0.81 at only 10\% recall. The model therefore operates as a
high-precision, low-recall filter at strict cutoffs: the top-ranked candidates are enriched for
true stabilizers, but the majority of stabilizing mutations are not recovered. We report this low
recall explicitly rather than obscuring it, because it defines the model's intended use as a
candidate-prioritization filter.

\subsection{Cross-stage comparison}
Stage~1 achieved a higher independent ROC AUC (0.732) than Stage~2 (0.669). However, the two
stages were trained and evaluated on different datasets, prediction targets (temperature
stability vs.\ folding \ddg), class labels, and sample sizes (1{,}563 vs.\ 669 records). The
comparison is therefore descriptive only: it is not a controlled head-to-head evaluation, and no
statistical significance test was performed between the two stages.
\fillerblock{confirm: original BRENDA size (inferred 2{,}446); positive/negative/ambiguous class definitions; true-\Tm{} subset $n$ for $r=0.60$; \ddg{} sign convention; how each Stage~2 threshold is applied; that thresholds were not selected on the benchmarks.}

% =====================================================================
\section{Discussion}
The two-stage framework provides a moderate, leakage-controlled route from whole-protein
screening to mutation-level optimization. Stage~1's independent ROC AUC of 0.732 indicates useful
but imperfect discrimination of thermostable proteins, and its precision--recall trade-off means
the operating point should be chosen to match the intended use (broad recall for discovery vs.\
high precision for shortlisting).

Stage~2's independent performance (Pearson 0.390; RMSE 1.509~kcal\,mol$^{-1}$; ROC AUC 0.669) is
best described as moderate, with pronounced high-precision/low-recall behaviour at strict \ddg{}
thresholds. This makes the model a candidate-prioritization filter rather than a comprehensive
predictor of every stabilizing mutation.

The central methodological contribution is explicit leakage auditing. Removing approximately
132{,}000 records---the great majority (131{,}479) detectable only through sequence homology---
shows that homology-level contamination, not duplicate records, is the dominant risk to honest
benchmarking here. A residual homology risk remains and is acknowledged rather than dismissed.

Comparisons with prior \Tm{} or \ddg{} models are only meaningful when the benchmark, sign
convention, data split, and metric are matched \filler{ref(s) for prior models, only if fairly comparable}. We therefore restrict quantitative comparison to genuinely comparable settings and
avoid implying superiority across incomparable evaluations.

We further caution that benchmark precision may not transfer unchanged to PET-degrading enzymes,
which are under-represented in these general stability sets. In particular, a benchmark precision
value such as 0.93 should not be equated with a 93\% experimental success rate; it is measured at
a chosen threshold on curated data, not in the laboratory.

Limitations include the moderate correlation on S669, the low recall at high-precision operating
points, the sequence-based (rather than structure-based) scope, residual homology risk, and the
fact that both benchmarks are general stability sets rather than PET-hydrolase-specific assays.

Future work should include PET-enzyme-specific evaluation, incorporation of the pH-optimum
analysis once a confirmed training run is available, and, ultimately, experimental validation
before any deployment claim.

% =====================================================================
\section{Methods}

\subsection{Study design}
The framework comprises two stages that use different prediction targets and different data
subsets: Stage~1 (whole-protein \Tm) and Stage~2 (mutation-level \ddg). The full 1{,}001{,}888-record
table did not train both models; records were routed by measurement type (Table~\ref{tab:data},
Panel~B).

\subsection{Source datasets}
Records were drawn from seven public sources spanning both measurement types
(Table~\ref{tab:data}, Panel~A), together with the independent BRENDA and S669 benchmarks and
UniProt for sequence resolution.
\fillerblock{for each of Domainome; Tsuboyama mega-scale; Tsuboyama doubles; Meltome Atlas; ThermoMutDB; FireProtDB; ProDDG/S2648; BRENDA; S669; UniProt --- give name, primary publication, repository, version/access date, measurement type, row count, inclusion/exclusion criteria, license, and DOI/accession.}

\subsection{Harmonization and sequence resolution}
Measurement typing, duplicate/replicate handling, and canonical-sequence resolution from UniProt
were applied to build the staging table.
\fillerblock{UniProt release; canonical vs.\ isoform handling; number and treatment of failed mappings; duplicate/replicate policy.}

\subsection{Leakage auditing}
Auditing combined an exact-key and mutation-level overlap check with a sequence-homology search;
records matching test proteins were removed, with counts recorded at each step.
\fillerblock{MMseqs2 version; identity threshold; coverage threshold; e-value; exact removal logic; counts removed/remaining at each step (total ~132{,}000; 131{,}479 by homology) --- taken from code.}

\subsection{Targets, features, models, and evaluation}
Stage~1 and Stage~2 targets were constructed from the BRENDA class definitions and the \ddg{}
mathematical definition, respectively, with the \degC{} and kcal\,mol$^{-1}$ thresholds applied as
reported in the Results.
\fillerblock{exact current feature set (if ESM-2 is used: model/version/layer/pooling/dimensionality/mutation representation); each model's package version, hyperparameters, class weighting, early stopping, scaling, and seed; ensemble type and weighting; validation/test split/grouping/fold/seed; metric definitions; bootstrap resampling unit, CI type, and seed; software/environment/hardware; ethics statement (computational, public data); substantive LLM-use disclosure per Nature Portfolio policy. Verify all from the current code; do not carry features from the 19,071-mutation model.}

% =====================================================================
\section*{Data Availability}
The processed training table, leakage-audited benchmark datasets, saved predictions, split
identifiers, leakage-audit outputs, and trained model artifacts generated in this study are
available through Zenodo (\url{https://doi.org/10.5281/zenodo.21257369}); source datasets are
cited in Methods and References. Deposit record: ``PET training set'' (Iyer, Ayush), one file,
approximately 2.5~GB.
\fillerblock{before submission: the Zenodo record currently requires login (restricted) --- make it open or provide a reviewer access mechanism; the deposit is currently GPL-3.0 (a software license) --- apply a data license such as CC~BY~4.0; confirm the single file contains all listed artifacts with a data dictionary/README; verify license and redistribution rights for each source dataset; remove EpHod unless the pH analysis is included.}

\section*{Code Availability}
Source code for data assembly, leakage auditing, model training, evaluation, threshold analysis,
bootstrap confidence-interval estimation, and figure generation is available at
\url{https://github.com/AyushIyer31/PET-Lab} and archived on Zenodo at
\url{https://doi.org/10.5281/zenodo.21519961} (release v1.0.0). Key scripts include
\texttt{ensemble\_eval.py}, \texttt{s669\_sweep.py}, \texttt{predict\_and\_sweep.py}, and
\texttt{bootstrap\_ci.py}.
\fillerblock{confirm archived release v1.0.0 (commit 497c4f5) is the exact code that produced the reported results; confirm each named script exists, matches the analysis, has no broken paths or exposed credentials, and runs with documented dependencies (README, license, environment file, seeds, hyperparameters, reproduction commands).}

% =====================================================================
\section*{Acknowledgements}
\fillerblock{brief; non-author contributors, technical/computational-resource support, and verified funding only. Only name Santa Clara University / UC Santa Cruz computing (e.g., the Nautilus cluster) or mentors with the exact program/allocation, named individuals, and their consent; do not imply endorsement.}

\section*{Author Contributions}
\fillerblock{provide CRediT roles for each author (choose from: Conceptualization; Data curation; Formal analysis; Funding acquisition; Investigation; Methodology; Project administration; Resources; Software; Supervision; Validation; Visualization; Writing --- original draft; Writing --- review \& editing). Use distinct forms for shared initials (Ayush Iyer / Abhinav Iyer / J.P.). Confirm author order.}
All authors reviewed and approved the final manuscript.

\section*{Additional Information}
\subsection*{Competing interests}
The authors declare no competing interests.
\fillerblock{all authors must confirm financial and non-financial competing interests, including any interest connected with PET-Lab; do not assume ``none''.}

% =====================================================================
\section*{References}
\fillerblock{References in numerical Nature style, cited sequentially, one source per number; every citation matched and every reference cited; no references in the Abstract. Replace each yellow ``[TO ADD: ref]'' marker above with a numbered citation. Include formal dataset citations for all seven sources and for the Zenodo data and code deposits. Verify authors, title, venue, year, volume, pages/article number, and DOI for every reference; do not fabricate DOIs.}

% Example scaffold (delete and replace with a real bibliography or \bibitem list):
\begin{enumerate}
  \item \filler{plastic-waste primary source (OECD/UNEP/peer-reviewed)}
  \item \filler{PETase discovery / engineering reference}
  \item \filler{Tsuboyama et al. 2023 mega-scale stability dataset}
  \item \filler{Meltome Atlas; ThermoMutDB; FireProtDB; ProDDG/S2648; Domainome; BRENDA; S669; UniProt (one entry each)}
  \item \filler{Zenodo data deposit (10.5281/zenodo.21257369) and code deposit (10.5281/zenodo.21519961)}
\end{enumerate}

% =====================================================================
% ============================ DISPLAY ITEMS ==========================
% =====================================================================
\clearpage
\section*{Tables}

% -------------------- TABLE 1 --------------------
\begin{table}[htbp]
\centering
\caption{Composition of the assembled dataset and allocation by measurement type.
Exact record counts; both panels sum to 1{,}001{,}888.}
\label{tab:data}
\small
\begin{tabular}{lr}
\toprule
\multicolumn{2}{l}{\textit{Panel A --- Rows by source dataset}}\\
\midrule
\textbf{Source dataset} & \textbf{Rows}\\
\midrule
Domainome (aPCA abundance)              & 457,943\\
Tsuboyama et al.\ 2023 (mega-scale)     & 357,155\\
Tsuboyama et al.\ 2023 (double mutants) & 138,275\\
Meltome Atlas                           & 27,884\\
ThermoMutDB                             & 11,520\\
FireProtDB                              & 6,549\\
ProDDG/S2648                            & 2,562\\
\midrule
\textbf{Total}                          & \textbf{1,001,888}\\
\bottomrule
\end{tabular}

\vspace{6pt}
\begin{tabular}{lrl}
\toprule
\multicolumn{3}{l}{\textit{Panel B --- Rows by measurement type and model allocation}}\\
\midrule
\textbf{Measurement type} & \textbf{Rows} & \textbf{Used to train}\\
\midrule
\ddg{} (mutation)     & 508,693 & Stage 2\\
Abundance (proxy)     & 457,943 & Not used in this study\\
\Tm{} (whole protein) & 29,654  & Stage 1\\
\dtm{}                & 5,598   & Not used in this study\\
\midrule
\textbf{Total}        & \textbf{1,001,888} & ---\\
\bottomrule
\end{tabular}

\vspace{3pt}
{\footnotesize \filler{confirm dataset versions and redistribution conditions per source}}
\end{table}

% -------------------- TABLE 2 --------------------
\begin{table}[htbp]
\centering
\caption{Stage 1 melting-temperature model --- classification on the independent BRENDA
benchmark (1{,}563 enzymes; 979 positive, 584 negative). ROC AUC $=0.732$ (95\% CI 0.708--0.755).
Positive class $=$ thermostable. Precision is reported with recall, F1, and accuracy so it is not
read as overall reliability.}
\label{tab:stage1}
\small
\begin{tabular}{lcccc}
\toprule
\multicolumn{5}{l}{\textit{Panel A --- Decision-threshold analysis}}\\
\midrule
\textbf{Threshold} & \textbf{Precision} & \textbf{Recall} & \textbf{F1} & \textbf{Accuracy}\\
\midrule
46 \degC & 0.65 & 0.98 & 0.78 & 0.66\\
50 \degC & 0.71 & 0.77 & 0.74 & 0.66\\
52 \degC & 0.77 & 0.63 & 0.69 & 0.65\\
60 \degC & 0.93 & 0.39 & 0.55 & 0.60\\
66 \degC & 0.98 & 0.32 & 0.49 & 0.58\\
\bottomrule
\end{tabular}

\vspace{6pt}
\begin{tabular}{lccc}
\toprule
\multicolumn{4}{l}{\textit{Panel B --- Per-model performance at the 60 \degC{} threshold}}\\
\midrule
\textbf{Model} & \textbf{Precision} & \textbf{Recall} & \textbf{F1}\\
\midrule
Extra Trees   & 0.963 & 0.369 & 0.533\\
Random Forest & 0.948 & 0.355 & 0.517\\
LightGBM      & 0.933 & 0.395 & 0.555\\
CatBoost      & 0.927 & 0.400 & 0.559\\
XGBoost       & 0.917 & 0.394 & 0.551\\
Ensemble      & 0.945 & 0.388 & 0.550\\
\bottomrule
\end{tabular}

\vspace{3pt}
{\footnotesize \filler{confirm BRENDA class definitions and that thresholds were not selected on this benchmark}}
\end{table}

% -------------------- TABLE 3 --------------------
\begin{table}[htbp]
\centering
\caption{Stage 2 mutation-level \ddg{} model --- regression and classification on the independent
S669 benchmark (669 mutations; 168 stabilizing, 501 non-stabilizing reference cases). Ensemble
ROC AUC $=0.669$ (95\% CI 0.62--0.72); ensemble regression Pearson $=0.390$ (95\% CI 0.32--0.46);
2{,}000 bootstrap resamples.}
\label{tab:stage2}
\small
\begin{tabular}{lccc}
\toprule
\multicolumn{4}{l}{\textit{Panel A --- Regression (predicted vs.\ experimental \ddg)}}\\
\midrule
\textbf{Model} & \textbf{Pearson} & \textbf{Spearman} & \textbf{RMSE (kcal\,mol$^{-1}$)}\\
\midrule
Extra Trees           & 0.436 & 0.461 & 1.481\\
XGBoost               & 0.382 & 0.413 & 1.514\\
Random Forest         & 0.369 & 0.396 & 1.523\\
LightGBM              & 0.330 & 0.364 & 1.549\\
CatBoost              & 0.324 & 0.362 & 1.557\\
Multilayer perceptron & 0.298 & 0.338 & 1.612\\
Ensemble              & 0.390 & 0.414 & 1.509\\
\bottomrule
\end{tabular}

\vspace{6pt}
\begin{tabular}{lcccc}
\toprule
\multicolumn{5}{l}{\textit{Panel B --- Ensemble classification across decision thresholds}}\\
\midrule
\textbf{Threshold} & \textbf{Precision} & \textbf{Recall} & \textbf{F1} & \textbf{Accuracy}\\
\midrule
0.00 kcal\,mol$^{-1}$  & 0.81 & 0.10 & 0.18 & 0.77\\
$+$0.25 kcal\,mol$^{-1}$ & 0.51 & 0.26 & 0.35 & 0.75\\
$+$0.50 kcal\,mol$^{-1}$ & 0.38 & 0.46 & 0.42 & 0.67\\
$+$1.00 kcal\,mol$^{-1}$ & 0.31 & 0.80 & 0.45 & 0.51\\
\bottomrule
\end{tabular}

\vspace{6pt}
\begin{tabular}{lcc}
\toprule
\multicolumn{3}{l}{\textit{Panel C --- Per-model performance at the strict (0.00 kcal\,mol$^{-1}$) threshold}}\\
\midrule
\textbf{Model} & \textbf{Precision} & \textbf{Recall}\\
\midrule
Random Forest         & 0.900 & 0.107\\
Extra Trees           & 0.786 & 0.065\\
XGBoost               & 0.735 & 0.149\\
LightGBM              & 0.625 & 0.119\\
CatBoost              & 0.429 & 0.107\\
Multilayer perceptron & 0.415 & 0.202\\
Ensemble              & 0.810 & 0.101\\
\bottomrule
\end{tabular}

\vspace{3pt}
{\footnotesize \filler{confirm \ddg{} sign convention and the predicted-stabilizer rule}}
\end{table}

% -------------------- FIGURE 1 --------------------
\clearpage
\section*{Figures}
\begin{figure}[htbp]
\centering
% Upload precision_vs_threshold.pdf to the Overleaf project root.
% (If it is not uploaded yet, a placeholder box is shown so the document still compiles.)
\IfFileExists{precision_vs_threshold.pdf}{%
  \includegraphics[width=\linewidth]{precision_vs_threshold.pdf}%
}{%
  \fbox{\parbox[c][0.32\linewidth][c]{0.96\linewidth}{\centering
    \filler{upload precision\_vs\_threshold.pdf to the Overleaf project root}}}%
}
\caption{\textbf{Precision as a function of the decision threshold for the two-stage framework on
independent benchmarks.} Precision for the Stage~1 whole-protein melting-temperature model on the
leakage-audited BRENDA benchmark ($n=1{,}563$; 979 positive, 584 negative) across
melting-temperature thresholds (panel~a) and for the Stage~2 mutation-level \ddg{} model on S669
($n=669$; 168 stabilizing, 501 non-stabilizing reference cases) across \ddg{} thresholds
(panel~b). In both panels the horizontal axis is oriented so that precision increases with more
stringent thresholds. Recall, F1, accuracy, and ROC AUC for the same evaluations are reported in
Tables~\ref{tab:stage1}--\ref{tab:stage2}. Positive class: BRENDA $=$ thermostable; S669 $=$
stabilizing.}
\label{fig:precision}
\end{figure}
\noindent{\footnotesize \filler{regenerate Figure 1 as a vector file from saved per-point predictions before submission}}

\end{document}
